% Irodov 第2.6节原题195--204，底本PDF页94--95，印刷页85--86。
\thermophaseproblem{195}{85}
假设水分子间的相互吸引作用突然消失，
其余宏观状态变量保持不变，
那么水施加到容器壁上的压强会增加多少？
\exsolution
范德瓦耳斯形式的状态方程可写为
\[
p=p_{\rm rep}-\frac{a}{V_m^2},
\]
其中 $V_m$ 为摩尔体积，
$p_{\rm rep}=RT/(V_m-b)$ 为排除体积修正的
有效斥力压强，
$a/V_m^2$ 代表分子间吸引的压强修正。
若吸引项消失而体积、温度瞬时保持不变，
墙压增加
\[
\boxed{\Delta p=\frac a{V_m^2}
=a\left(\frac{\rho}{M}\right)^2}.
\]
对密度约
$\rho\simeq1.0\,\mathrm{kg/L}$ 的水，
摩尔质量 $M\simeq18\,\mathrm{g/mol}$，
取水的范德瓦耳斯参数
$a\simeq5.5\,\mathrm{atm\,L^2/mol^2}$，
则
$V_m\simeq0.018\,\mathrm{L/mol}$，
\[
\boxed{\Delta p\simeq1.7\times10^4\,\mathrm{atm}}.
\]
\excheck 这是范德瓦耳斯模型的瞬态分子作用
\emph{形式估算}；
现实水移除吸引作用后通常不可能维持原有液态结构。

\thermophaseproblem{196}{85}
液体密度为 $\rho$，
单位质量汽化潜热为 $q$。
假设汽化所需潜热主要用于克服液体的
内聚压强 $p_i$，
且液体服从范德瓦耳斯模型，
求 $p_i$ 的近似值，并估算水的数值。
\exsolution
取质量为 $m$ 的液体，
汽化到密度远小于液体的蒸汽时，
克服内聚压力的机械能约为
$p_i$ 乘初始液体的相应有效体积
$m/\rho$。
按题设将潜热近似视为内聚作用功，
\[
mq\simeq p_i\frac m\rho.
\]
消去质量：
\[
\boxed{p_i\simeq\rho q}.
\]
对于水，
$\rho\simeq10^3\,\mathrm{kg/m^3}$、
$q\simeq2.26\times10^6\,\mathrm{J/kg}$，
故
\[
\boxed{p_i\simeq2.3\times10^9\,\mathrm{Pa}
\simeq2.2\times10^4\,\mathrm{atm}}.
\]
\excheck 题给将潜热全归于内聚功属于近似假定，
真实汽化焓还包含对外做功和分子内能等贡献。

\thermophaseproblem{197}{85}
证明：一摩尔范德瓦耳斯气体在临界状态满足
\[
V_{m,c}=3b,\qquad
p_c=\frac{a}{27b^2},\qquad
T_c=\frac{8a}{27Rb},
\qquad
p_cV_{m,c}=\frac38RT_c.
\]
\exsolution
一摩尔范德瓦耳斯状态方程：
\[
p(V_m,T)=\frac{RT}{V_m-b}-\frac{a}{V_m^2}.
\]
临界等温线出现水平拐点，故
\[
\left(\frac{\partial p}{\partial V_m}\right)_{T_c}
=-\frac{RT_c}{(V_m-b)^2}+\frac{2a}{V_m^3}=0,
\]
\[
\left(\frac{\partial^2p}{\partial V_m^2}\right)_{T_c}
=\frac{2RT_c}{(V_m-b)^3}-\frac{6a}{V_m^4}=0.
\]
由第一式
$RT_c=2a(V_m-b)^2/V_m^3$，
由第二式
$RT_c=3a(V_m-b)^3/V_m^4$。
两式相等，约去公共因子：
$2V_m=3(V_m-b)$，
所以
\[
\boxed{V_{m,c}=3b}.
\]
代入第一式得
$RT_c=8a/(27b)$，
再代入状态方程得
$p_c=a/(27b^2)$。
于是
\[
\boxed{p_cV_{m,c}
=\frac a{27b^2}(3b)
=\frac a{9b}
=\frac38RT_c}.
\]
\excheck 临界点是等温线的\emph{水平拐点}，
因此必须同时使用一阶、二阶导数为零，
不能仅仅令斜率为零。

\thermophaseproblem{198}{85}
已知二氧化碳的临界温度
$T_c=304\,\mathrm K$、
临界压强 $p_c=73\,\mathrm{atm}$。
求范德瓦耳斯常数 $a,b$。
\exsolution
由上一题临界关系，
$p_c=a/(27b^2)$、
$RT_c=8a/(27b)$。
先由两式联立消去 $a$：
\[
\boxed{b=\frac{RT_c}{8p_c}},
\]
再由 $a=27p_cb^2$：
\[
\boxed{a=\frac{27R^2T_c^2}{64p_c}}.
\]
以
$R=0.082057\,\mathrm{L\,atm/(mol\,K)}$
计算：
\[
\boxed{b\simeq0.0427\,\mathrm{L/mol}},
\qquad
\boxed{a\simeq3.6\,\mathrm{atm\,L^2/mol^2}}.
\]
\excheck $a$ 的量纲为压强乘摩尔体积的平方；
$b$ 的量纲为摩尔体积。

\thermophaseproblem{199}{85}
苯 $\mathrm{C_6H_6}$ 的临界温度
$T_c=562\,\mathrm K$、
临界压强 $p_c=47\,\mathrm{atm}$。
若按范德瓦耳斯模型，
求苯临界态的质量比体积 $v_c$。
\exsolution
由临界恒等式
$p_cV_{m,c}=(3/8)RT_c$，
临界摩尔体积为
$V_{m,c}=3RT_c/(8p_c)$。
质量比体积是单位质量物质的体积，
\[
v_c=\frac{V_{m,c}}M
=\boxed{\frac{3RT_c}{8Mp_c}}.
\]
取苯摩尔质量
$M\simeq78\,\mathrm{g/mol}$、
$R=0.082057\,\mathrm{L\,atm/(mol\,K)}$，
得到
\[
\boxed{v_c\simeq4.7\times10^{-3}\,\mathrm{L/g}
=4.7\,\mathrm{cm^3/g}}.
\]
\excheck 临界\emph{质量比体积}
与摩尔体积不同，必须除以摩尔质量 $M$。

\thermophaseproblem{200}{85}
将范德瓦耳斯状态方程改写为约化变量形式，
其中
$\pi=p/p_c$、
$\upsilon=V_m/V_{m,c}$、
$\tau=T/T_c$。
如果气体压强达到临界压强的12倍，
而体积是临界体积的一半，
求温度与临界温度之比。
\exsolution
原方程
$(p+a/V_m^2)(V_m-b)=RT$。
代入临界关系
$V_{m,c}=3b$、
$p_c=a/(27b^2)$、
$RT_c=8a/(27b)$：
\[
p=p_c\pi,\quad V_m=3b\upsilon,\quad T=T_c\tau.
\]
得
\[
\left[p_c\pi+\frac a{9b^2\upsilon^2}\right]
(3b\upsilon-b)=RT_c\tau.
\]
用 $a/(9b^2)=3p_c$、
$RT_c=8p_cb$，化为
\[
\boxed{\left(\pi+\frac3{\upsilon^2}\right)
(3\upsilon-1)=8\tau}.
\]
给定 $\pi=12$、$\upsilon=1/2$，
\[
8\tau=(12+3/(1/2)^2)(3/2-1)
=(12+12)\frac12=12,
\]
故
\[
\boxed{\tau=\frac32=1.5}.
\]
\excheck $\upsilon=1/2>1/3$，
满足范德瓦耳斯方程 $V_m>b$ 的必要体积约束。

\thermophaseproblem{201}{85}
已知水的范德瓦耳斯常数 $a,b$。
求：(a) 质量
$m=1.00\,\mathrm{kg}$ 的水在仍作为液体时
可占据的最大体积；
(b) 饱和水蒸气压强的最大可能值
（范德瓦耳斯模型临界上界）。
\exsolution
范德瓦耳斯模型的液气共存支路在临界点终止。
气、液态临界摩尔体积为
$V_{m,c}=3b$。
对总物质的量
$\nu=m/M$，临界总体积
\[
\boxed{V_{\max,\ell}=\nu(3b)=\frac{3bm}{M}}.
\]
当温度达到临界值时，
饱和蒸汽压强达到临界压强：
\[
\boxed{p_{\rm sat,max}=p_c=\frac a{27b^2}}.
\]
以水在模型中的常用近似参数
$b\simeq0.030\,\mathrm{L/mol}$、
$a\simeq5.5\,\mathrm{atm\,L^2/mol^2}$，
$M=0.018\,\mathrm{kg/mol}$，
\[
\boxed{V_{\max,\ell}\simeq5.0\,\mathrm L},\qquad
\boxed{p_{\rm sat,max}\simeq2.3\times10^2\,\mathrm{atm}}.
\]
\excheck 这里的“最大液体体积”指
范德瓦耳斯\emph{液相支路延至临界点}，
不是任意温度或压力下液态水体积上限。

\thermophaseproblem{202}{85}
把二氧化碳视为范德瓦耳斯气体，
利用其参数 $a,b$ 求临界温度和临界质量密度。
\exsolution
对每摩尔范德瓦耳斯气体，
\[
T_c=\frac{8a}{27Rb},\qquad V_{m,c}=3b.
\]
临界质量密度为
\[
\boxed{\rho_c=\frac M{V_{m,c}}=\frac M{3b}}.
\]
取二氧化碳
$a\simeq3.59\,\mathrm{atm\,L^2/mol^2}$、
$b\simeq0.0427\,\mathrm{L/mol}$、
$M\simeq44.0\,\mathrm{g/mol}$，
以适配的 $R=0.082057$，
\[
\boxed{T_c\simeq3.0\times10^2\,\mathrm K
=0.30\,\mathrm{kK}},
\]
\[
\boxed{\rho_c\simeq\frac{44}{3(0.0427)}
\,\mathrm{g/L}\simeq0.34\,\mathrm{g/cm^3}}.
\]
\excheck 摩尔质量与 $b$ 均须使用同一物质，
单位换算
$1000\,\mathrm{g/L}=1\,\mathrm{g/cm^3}$。

\thermophaseproblem{203}{85}
室温下盛有一部分液态乙醚的密闭容器，
准备将其加热至临界状态。
乙醚临界温度
$T_c=467\,\mathrm K$、
临界压强
$p_c=35.5\,\mathrm{atm}$、
摩尔质量
$M=74\,\mathrm{g/mol}$。
求室温时液态乙醚所占容器体积的比例。
\exsolution
设容器总体积为 $V_{\rm vessel}$，
乙醚质量为 $m$，
若总平均比容需要等于临界比容，
则
$V_{\rm vessel}=mV_{m,c}/M$。
而室温液态乙醚自身占据体积
$V_\ell=m/\rho_\ell$。
所以所求体积分数
\[
f=\frac{V_\ell}{V_{\rm vessel}}
=\frac M{\rho_\ell V_{m,c}}.
\]
范德瓦耳斯临界关系
$V_{m,c}=3RT_c/(8p_c)$，
所以
\[
\boxed{f=\frac{8Mp_c}{3\rho_\ell RT_c}}.
\]
以室温乙醚密度
$\rho_\ell\simeq0.71\,\mathrm{g/cm^3}
=710\,\mathrm{g/L}$ 计算，
\[
\boxed{f\simeq0.25}.
\]
\excheck 乙醚室温密度\emph{并未在题面中给出}，
因此数值需要外部物性表或教材常用近似；
公式本身只依赖临界关系和 $\rho_\ell$。

\thermophaseproblem{204}{85}
证明：图2.5中连接气液两相共存态1与5的
等温等压水平线，
会把范德瓦耳斯等温曲线夹出的
区域 I 与 II 分为相同面积
（Maxwell 等面积构造）。
\exsolution
设共存温度为 $T$、
共存压强为 $p_s$，
液相端摩尔体积为 $V_\ell$，
气相端摩尔体积为 $V_g$。
对于沿平衡态、可逆等温路径的
Helmholtz 自由能有
\[
\mathrm dF=-p\,\mathrm dV
\quad(T=\text{常数}).
\]
共存两相在相同 $T,p_s$ 下，
其摩尔 Gibbs 自由能相同：
\[
F(V_g)+p_sV_g=F(V_\ell)+p_sV_\ell.
\]
故
\[
F(V_g)-F(V_\ell)
=-p_s(V_g-V_\ell).
\]
另一方面沿范德瓦耳斯等温曲线
形式积分：
\[
F(V_g)-F(V_\ell)
=-\int_{V_\ell}^{V_g}p_{\rm vdW}(V,T)\,\mathrm dV.
\]
将两式相等：
\[
\boxed{\int_{V_\ell}^{V_g}
[p_{\rm vdW}(V,T)-p_s]\,\mathrm dV=0}.
\]
这意味着等温曲线位于水平线之上的
有符号面积与位于其下方的面积模相等，
即图示区域 I 与 II 面积相等。
\excheck 等面积规则是由\emph{共存端点的
自由能相等}导出；
不能把范德瓦耳斯中间不稳定分支
误当作真实液气共存时可以稳定走过的物理过程。
